\documentclass[11pt]{article}

\usepackage[margin=1in]{geometry}
\usepackage{amsmath, amssymb}
\usepackage{array}
\usepackage{enumitem}
\usepackage{hyperref}
\hypersetup{hidelinks}

\title{Scaling Book Exercises: Section 2}
\author{Lucas Yan}
\date{}

\newcommand{\tmath}{T_{\mathrm{math}}}
\newcommand{\tcomm}{T_{\mathrm{comm}}}
\newcommand{\tmem}{T_{\mathrm{mem}}}

\begin{document}
\maketitle

\noindent
\textbf{Source scan:} \texttt{../handwritten/section\_02\_handwritten.pdf}

\section*{Exercise 1}

For v4p HBM bandwidth,
\[
\mathrm{BW}_{\mathrm{HBM}} \approx 1.2 \times 10^{12} \text{ bytes/s}.
\]
Moving $200 \mathrm{GB}$ in two directions gives
\[
200 \mathrm{GB} \times 2 = 400 \times 10^9 \text{ bytes}.
\]
Across $32$ chips, this is
\[
\frac{400 \times 10^9}{32} = 12.5 \times 10^9
\]
bytes per chip.
Therefore,
\[
\text{total time}
  =
  \frac{12.5 \times 10^9}{1.2 \times 10^{12}}
  =
  \approx 10^{-2} \text{ s}
  \approx 10 \text{ ms}.
\]

\section*{Exercise 2}

The scanned notes compare v5e and v5p hardware.

\[
\begin{array}{c|cc}
 & \mathrm{v5e} & \mathrm{v5p} \\
\hline
\mathrm{cores} & 32 & 2240 \\
\mathrm{TPU\ tensor\ cores} & 256 & 17920 \\
\mathrm{total\ FLOPs/s} & 5.1 \times 10^{16} & 4.1 \times 10^{18} \\
\mathrm{total\ HBM} & 4\ \mathrm{TB} & 860\ \mathrm{TB}
\end{array}
\]

For v5e, one calculation in the notes is
\[
\frac{16 \cdot 16}{4 \cdot 2} = 32.
\]

For v5p, the notes use:
\begin{align*}
\text{(a)}\quad & \frac{16 \cdot 20 \cdot 28}{2 \cdot 2 \cdot 1}
  = 4 \cdot 20 \cdot 28
  = 2240,\\
\text{(b)}\quad & 2 \cdot 16 \cdot 20 \cdot 28
  = 1120 \cdot 16
  = 17920,\\
\text{(c)}\quad & 16 \cdot 20 \cdot 28 \cdot 4.59 \times 10^{14}
  \approx 4.0 \times 10^{18},\\
\text{(d)}\quad & 96 \cdot 16 \cdot 20 \cdot 28
  = 860160\ \mathrm{GB}
  \approx 860\ \mathrm{TB}.
\end{align*}

\section*{Exercise 3}

The computation time is
\[
\mathrm{FLOPs} = 2BDF,
\qquad
\mathrm{FLOPs/s}_{\mathrm{v6e}} = 9.2 \times 10^{14},
\]
so, for the scanned setting,
\[
\tmath = \frac{2BDF}{9.2 \times 10^{14}}
       = \frac{8BD^2}{9.2 \times 10^{14}}.
\]

For memory communication, the notes use
\[
2DF + 2BD + 2BF = 2(BD + DF + BF)
\]
bytes and PCIe bandwidth
\[
1.6 \times 10^{10} \text{ bytes/s}.
\]
Thus,
\[
\tcomm
  =
  \frac{2(BD + DF + BF)}{1.6 \times 10^{10}}.
\]

The compute-bound condition is
\[
\tmath > \tcomm.
\]
Using $F = 4D$ and $B \ll D$, this becomes
\[
\frac{8BD^2}{9.2 \times 10^{14}}
  >
\frac{8D^2}{1.6 \times 10^{10}},
\]
and therefore
\[
B > \frac{9.2 \times 10^{14}}{16 \times 10^9}
  = 5.75 \times 10^4
  \approx 57500.
\]

\section*{Exercise 4}

\subsection*{Part (a)}

For a matrix multiplication with dimensions shown in the scan,
\[
\mathrm{FLOPs}
  =
  2 \cdot B \cdot 4096 \cdot 16384
  =
  134217728 B
  \approx 1.3 \times 10^8 B.
\]

Using v5e throughput,
\[
\mathrm{FLOPs/s}_{\mathrm{v5e}} = 3.94 \times 10^{14},
\]
the compute time is
\[
\tmath = \frac{1.3 \times 10^8}{3.94 \times 10^{14}} B.
\]

The memory time is modeled as
\[
\tmem
  =
  \frac{16384 \cdot 4096 + 4096B + 16384B}{8.2 \times 10^{11}}
  \approx
  \frac{6.7 \times 10^7 + 20000B}{8.2 \times 10^{11}}.
\]

We are compute-bound if
\[
\tmath > \tmem.
\]
Solving the inequality as in the scan gives
\[
B > 267.
\]

\subsection*{Part (b)}

If loading is $22\times$ faster, then
\[
\tmem' = \frac{\tmem}{22}.
\]
Using the same inequality with the VMEM bandwidth denominator increased from
$8.2 \times 10^{11}$ to roughly $1.80 \times 10^{13}$, the compute-bound
threshold becomes
\[
B > 11.
\]
In practice, since all VMEM bandwidth cannot usually be dedicated to this one
load, the practical threshold is closer to $20$.

\section*{Exercise 5}

\begin{enumerate}[label=(\alph*)]
  \item With $3$ rows and $3$ columns, there are
  \[
  3 + 3 = 6
  \]
  links.

  \item The total bytes are
  \[
  2 \cdot 8 \cdot 128 \cdot 8192
  \approx 1.7 \times 10^7 \text{ bytes}.
  \]

  The diagram in the notes routes from $(0,0)$ through $(0,1)$ and then to $(1,0)$, so the data can be sent two ways. With total bandwidth
  \[
  2 \cdot 4.5 \times 10^{10} = 9.0 \times 10^{10},
  \]
  the transfer time is
  \[
  \frac{1.7 \times 10^7}{9.0 \times 10^{10}}
  \approx 188 \mu\mathrm{s}.
  \]
  The notes record a total time of approximately
  \[
  188 \mu\mathrm{s} + 6 \mu\mathrm{s} = 194 \mu\mathrm{s}.
  \]
\end{enumerate}

\section*{Exercise 6}

The data path in the notes is:
\[
\text{host DRAM}
\rightarrow
16 \text{ TPU HBM shards}
\rightarrow
\text{TPU }(0,0)\text{ HBM}
\rightarrow
\text{MXU matrix multiplication}.
\]

The array is approximately
\[
A \approx 17 \mathrm{GB},
\qquad
x = 2 \mathrm{MB}.
\]
Each of the $16$ TPUs holds approximately $1 \mathrm{GB}$.

\subsection*{(i) Host DRAM to TPU HBM shards}

Using PCIe bandwidth from Exercise 3,
\[
\mathrm{BW}_{\mathrm{PCIe}} = 1.6 \times 10^{10} \text{ bytes/s}
  = 16 \mathrm{GB/s}.
\]
With $16$ PCIe lanes, the total bandwidth is approximately
\[
16 \cdot 16 = 256 \mathrm{GB/s}.
\]
Thus,
\[
\text{time}_{\mathrm{PCIe}}
  =
  \frac{16}{256}
  =
  \frac{1}{16}
  =
  0.0625\text{ s}
  =
  62.5\text{ ms}.
\]

\subsection*{(ii) Receive remaining data from other TPUs}

The $(0,0)$ TPU must receive the remaining $15 \mathrm{GB}$ from other TPUs. Each link is approximately $45 \mathrm{GB/s}$, so two links provide
\[
90 \mathrm{GB/s}.
\]
Therefore,
\[
\text{time}
  =
  \frac{15}{90}
  =
  \frac{1}{6}
  \approx
  0.167\text{ s}
  =
  167\text{ ms}.
\]

\subsection*{(iii) HBM to MXU}

Moving $16 \mathrm{GB}$ from HBM to MXU with HBM bandwidth
\[
8.2 \times 10^{11} \text{ bytes/s}
\]
takes
\[
\text{time}
  =
  \frac{16 \times 10^9}{8.2 \times 10^{11}}
  \approx
  2 \times 10^{-2}\text{ s}
  =
  20\text{ ms}.
\]

\subsection*{(iv) Compute}

The compute required is
\[
2 \cdot 8 \cdot (128 \cdot 1024)^2
  =
  2.75 \times 10^{11}\text{ FLOPs}.
\]
Using throughput
\[
1.97 \times 10^{14}\text{ FLOPs/s},
\]
the compute time is
\[
\frac{2.75 \times 10^{11}}{1.97 \times 10^{14}}
  =
  1.4 \times 10^{-3}\text{ s}
  =
  1.4\text{ ms}.
\]

The component times are
\[
62.5\text{ ms}, \qquad
167\text{ ms}, \qquad
20\text{ ms}, \qquad
1.4\text{ ms}.
\]
Assuming perfect overlap and pipelining, the lower bound is
\[
\max(62.5, 167, 20, 1.4)\text{ ms}
  =
  167\text{ ms}.
\]

\section*{Verification Notes}

The following items should be verified against the exercise text before final submission:
\begin{itemize}
  \item The exact HBM bandwidth value in Exercise 1; the official text rounds the v4p value to about $1.2 \times 10^{12}$ bytes/s and states a roughly $10$ ms result.
  \item Exercise 5's topology labels and whether the $6\mu$s latency term is per-hop or total.
\end{itemize}

\end{document}
